\section{Experiments}

In this section, we present the experimental setup and results of our approach. We evaluated our method only on instances of the \href{https://huggingface.co/datasets/neo4j/text2cypher-2024v1}{neo4j/text2cypher-2024v1} test set that have an already deployed Neo4j database available on \href{https://sandbox.neo4j.com/}{Neo4j Sandbox}. This reduced the number of test instances from 4,833 to 2471, but it was strictly necessary to, at least, execute our agentic approach against the actual database used for the generation of samples in the test set. Furthermore, we used already computed zero-shot and finetuned results reported in the original paper \cite{neo4jlabs-text2cypher-2024} as several baselines to compare our approach with their one.

\subsection{Evaluation Metrics}

We evaluate our text-to-Cypher approach using a combination of runtime-based and structural metrics, following the evaluation protocols established in prior work \cite{neo4jlabs-text2cypher-2024, cypherbench-megagonlabs-2024}. These metrics collectively assess both the syntactic correctness and semantic equivalence of generated Cypher queries.

\paragraph{Execution Accuracy (EA).} This metric measures whether a generated Cypher query executes successfully on the target Neo4j database without raising syntax or runtime errors. A query is considered valid if it can be executed within a specified timeout (120 seconds in our experiments) without encountering exceptions such as \texttt{CypherSyntaxError}, \texttt{DatabaseError}, \texttt{CypherTypeError}, or \texttt{ClientError}. Execution Accuracy provides a fundamental threshold for query quality, as non-executable queries offer no practical value regardless of their structural similarity to ground truth.

\paragraph{Exact Match (EM).} The Exact Match metric represents the most stringent evaluation criterion, requiring character-level equivalence between the predicted and ground truth Cypher queries. While highly precise, this metric is known to be overly conservative, as semantically equivalent queries may differ in syntactic representation (e.g., variable naming, clause ordering, or equivalent formulations). Nevertheless, EM serves as an important upper bound for perfect generation accuracy.

\paragraph{Result Similarity (RS).} To assess semantic equivalence at the data level, we compute the Jaccard similarity between the result sets returned by the predicted and ground truth queries. Specifically, we represent each query result as a list of dictionaries, where each dictionary corresponds to a row in the output. The similarity computation accounts for row ordering when either query contains an \texttt{ORDER BY} clause; otherwise, we apply a greedy alignment algorithm that pairs rows from both result sets to maximize overall similarity. For each aligned pair of rows, we convert values to hashable types and compute their Jaccard index. The final Result Similarity is the Jaccard coefficient over all (row index, value) pairs across both result sets. This metric captures whether queries retrieve the same data, even when their structural implementations differ. Following \cite{neo4jlabs-text2cypher-2024}, a value of RS $\geq 0.5$ is considered a successful match.

\paragraph{Provenance Subgraph Jaccard Similarity (PSJS).} Introduced by \cite{cypherbench-megagonlabs-2024}, this metric evaluates the structural overlap of the graph patterns matched by two queries, rather than their final result sets. The provenance subgraph of a query consists of all nodes and relationships traversed during pattern matching, regardless of what is ultimately returned. For each query, we extract the \texttt{MATCH} and \texttt{WHERE} clauses, augment unlabeled nodes and relationships with temporary variables, and construct a Cypher query that returns the element IDs of all matched nodes. The PSJS is then computed as the Jaccard coefficient between the sets of element IDs from the predicted and ground truth provenance subgraphs:
\[
\text{PSJS} = \frac{|PS_{\text{pred}} \cap PS_{\text{target}}|}{|PS_{\text{pred}} \cup PS_{\text{target}}|}
\]
where $PS_{\text{pred}}$ and $PS_{\text{target}}$ denote the provenance subgraphs of the predicted and target queries, respectively. This metric is particularly valuable for evaluating whether a model correctly understands the graph structure to be queried, even if the final projection or aggregation differs from the ground truth. As with Result Similarity, we adopt the threshold of PSJS $\geq 0.5$ as a success criterion following \cite{cypherbench-megagonlabs-2024}.

Together, these metrics provide a multi-faceted evaluation framework that captures syntactic validity (EA, EM), semantic data equivalence (RS), and structural query understanding (PSJS), enabling comprehensive assessment of text-to-Cypher generation quality.

\subsection{Results}